\mode*

\section{Disposition}

\begin{frame}
  \begin{figure}
    \includegraphics[height=0.8\textheight]{fig/book.jpg}
    \caption{Front cover of Necessary Conditions of Learning.}
  \end{figure}
\end{frame}

\begin{frame}
  \begin{enumerate}
    \alert{\item What makes humans human?}
    \item What is to be learned?
    \item Sameness and difference in learning
    \item What does the world look like to others?
    \item The art of learning
    \item Making learning possible
    \item Learning to help others to learn
  \end{enumerate}
\end{frame}

\begin{frame}
  \begin{block}{Ch 1 What makes humans human?}
    \begin{itemize}
      \item Pedagogy and advanced learning makes humans human.
      \item What is pedagogy?
      \item How does it relate to learning?
    \end{itemize}
  \end{block}
\end{frame}


\section{Pedagogy}

\begin{frame}
  \begin{block}{Cultural evolution}
    \begin{itemize}
      \item We pass on innovations from generation to generation.
      \item Done through pedagogy.
    \end{itemize}
  \end{block}
\end{frame}

\begin{frame}
  \begin{definition}[Premack's 1984 definition]
    \begin{enumerate}
      \item There is a goal for learning.
      \item There is a systematic attempt (by the \enquote{teacher}) to help 
        the learner to reach that goal.
      \item The systematic attempt is guided by the teacher's perception of the 
        learner's progress (or lack of progress).
    \end{enumerate}
  \end{definition}

  \pause

  \begin{remark}
    \begin{itemize}
      \item This is what makes us different from animals.
      \item Teaching and learning seems innate.
    \end{itemize}
  \end{remark}
\end{frame}

\subsection{Possibility of learning from others}

\begin{frame}
  \begin{example}[Learn with pedagogy]
    \begin{itemize}
      \item We've all passed through school.
      \item Someone has explained something for us.
      \item Or we've explained something to someone.
    \end{itemize}
  \end{example}

  \pause

  \begin{remark}
    \begin{itemize}
      \item We'll revisit this in Chapter 6--7.
    \end{itemize}
  \end{remark}
\end{frame}

\begin{frame}
  \begin{example}[Learn without pedagogy]
    \begin{itemize}
      \item We learn by research.
      \item Children learn by experimenting on their own.
    \end{itemize}
  \end{example}

  \pause
  
  \begin{remark}
    \begin{itemize}
      \item We'll revisit this in Chapter 5.
    \end{itemize}
  \end{remark}
\end{frame}

\subsection{Pedagogy or no pedagogy?}

\begin{frame}
  \begin{example}[Learning as a by-product or as an aim]
    \begin{itemize}
      \item Apprenticeship (PhD studies?) is learning as a by-product.
      \item As opposed to schooling, which is learning as an aim.
    \end{itemize}
  \end{example}

  \pause

  \begin{example}[\enquote{De-pedagogized} learning]
    \begin{itemize}
      \item Authentic vs unauthentic learning
      \item \enquote{Real life} vs school
      \item Problem/project-based learning?
    \end{itemize}
  \end{example}

  \pause

  \begin{question}
    \begin{itemize}
      \item Is one better than the other?
    \end{itemize}
  \end{question}
\end{frame}


\section{Pedagogies of learning}

\begin{frame}
  \begin{remark}[Institutionalization]
    \begin{itemize}
      \item Grades
      \item Selection of students
      \item School rankings
      \item \etc
    \end{itemize}
  \end{remark}
\end{frame}

\begin{frame}
  \begin{definition}[Double-sided paradox]
    \begin{itemize}
      \alert<1>{\item Teacher's paradox.}
      \alert<2>{\item Student's paradox}
    \end{itemize}
  \end{definition}

  \begin{onlyenv}<1>
    \begin{example}[The teacher's paradox]
      \blockcquote[p.~13]{NecessaryConditionsOfLearning}{%
        \textins{T}he more the teacher does to enable the students to answer 
        the questions asked, to solve the problems given, the fewer 
        opportunities may be left for the students to learn to understand that 
        which they are expected to learn to understand. The reason is that the 
        more clearly the teacher tells the students what is to be done, the 
        less chance the students get to make the necessary distinctions (for 
        instance, between what is critical and what is not).%
      }
    \end{example}
  \end{onlyenv}

  \begin{onlyenv}<2>
    \begin{example}[The student's paradox]
      \blockcquote[p.~13]{NecessaryConditionsOfLearning}{%
        \textins{T}hey are expected to learn (to understand), but the more the 
        teacher points out what they have to do or say, the less opportunity 
        they have to learn (to understand), as they have fewer distinctions to 
        make.%
      }
    \end{example}
  \end{onlyenv}
\end{frame}

\begin{frame}
  \begin{remark}[The problem]
    \begin{itemize}
      \item Surface vs deep learning~\cite{DeepSurfaceLearning}.
      \item Learn what to do vs learn how to see.
      \item \enquote{Teach for the test!}
    \end{itemize}
  \end{remark}
\end{frame}

\section{Conclusion}

\begin{frame}
  \begin{block}{In conclusion}
    \begin{itemize}
      \item Pedagogy is innate in our species.
      \item We can learn without pedagogy too.
      \item Pedagogy has advantages, but is not without problems.
      \item This book focuses on learning:
        \begin{itemize}
          \item Focus is on how to make our pedagogy yield the best learning.
          \item We'll have to revisit both.
        \end{itemize}
    \end{itemize}
  \end{block}
\end{frame}

\mode<all>{\endinput}

\begin{frame}
  \begin{onlyenv}<1>
    \begin{remark}
      \blockcquote[p.~21]{NecessaryConditionsOfLearning}{%
        Teachers must somehow organize work in the classroom in order to make 
        learning possible.
        A frequent question posed in educational research is:
        Which is the best arrangement for learning?%
      }
    \end{remark}
  \end{onlyenv}
  \begin{example}[{\cite[p.~21]{NecessaryConditionsOfLearning}}]
    \begin{itemize}
      \item Is learning by yourself better than learning by being taught?
      \item Does homework enhance learning?
      \item Is problem-based learning better than lectures for big classes?
      \item Is individualized learning preferable to group work?
      \item Is project work a good idea?
    \end{itemize}
  \end{example}
  \begin{onlyenv}<2>
    \begin{remark}
      \blockcquote[p.~21]{NecessaryConditionsOfLearning}{%
        The problem with questions of this kind is that they cannot be 
        answered.
        It is not that they cannot be answered \emph{yet}, and it is not 
        because of a scarcity of research funds or a scarcity of good ideas.
        They are simply imponderable because of the degree of generality.%
      }
    \end{remark}
  \end{onlyenv}
\end{frame}

\section{Take away}

\begin{frame}
  \begin{center}
    What is to be learned?\\[1em]
    instead of\\[1em]
    What is to be done?
  \end{center}
\end{frame}

\begin{frame}
  \begin{block}{Ch 2 What is to be learned?}
    \begin{itemize}
      \item What is an intended learning outcome?
      \item What should it be?
      \item It depends on the learner: \emph{critical aspects, critical 
        features} --- the atoms of the learning objective.
    \end{itemize}
  \end{block}
\end{frame}
